\documentclass[../main.tex]{subfiles}

\begin{document}

\chapter{Matrices}

Matrices are one of the fundamental components that build linear
algebra. This note discusses the fundamental concepts and properties
of matrices.

\section{Basic Terms and Notations}

As always, let's get started with the definition.

\begin{keydef}
A \textbf{matrix} is a 2-dimensional array composed with entries
called \textbf{elements}. A matrix is said to have \textbf{order}
``$m$ by $n$'' and written as $m \times n$ if there are $m$ rows
and $n$ columns.
\end{keydef}

Technically, the elements of a matrix can be anything including
functions. However, majority of the matrices discussed in
linear algebra will have numeric elements. Let's take a look
at some examples.
\[
A=
\begin{bNiceMatrix}[last-row,last-col]
    1  &  2 &  3 & 4  &
        \Vbrace[horizontal-label]{3}{3 \text{ rows}} \\
    -1 & -2 & -3 & -4 \\
    0  & -1 & 0  & 1  \\
    \Hbrace{4}{4 \text{ columns}} &
\end{bNiceMatrix}
\qquad\quad,\qquad\qquad
B =
\begin{bNiceMatrix}[last-row,last-col]
\pi & \Vbrace[horizontal-label]{4}{4 \text{ rows}} \\
e \\
0 \\
1 \\
\Hbrace{1}{1 \text{ column}}
\end{bNiceMatrix}
\]
Here, matrix $A$ has order $3 \times 4$ and matrix $B$ has
order $4 \times 1$. The from of a matrix that has order
$m \times n$ can be generalized as the following.
\[
A=
\begin{bNiceMatrix}[last-row,last-col]
    a_{1,1} & a_{1,2} & a_{1,3} & \cdots & a_{1,n}
        & \Vbrace[horizontal-label]{5}{n \text{ rows}} \\
    a_{2,1} & a_{2,2} & a_{2,3} & \cdots & a_{2,n} \\
    a_{3,1} & a_{3,2} & a_{3,3} & \cdots & a_{3,n} \\
    \vdots  & \vdots  & \vdots  & \ddots & \vdots  \\
    a_{m,1} & a_{m,2} & a_{m,3} & \cdots & a_{m,n} \\
    \Hbrace{5}{m \text{ columns}} &
\end{bNiceMatrix}
\]
It is also worth noting that matrix $A$ is often represented
as $[a_{ij}]$ and the elements as $a_{ij}$ like $a_{12}$ for
$a_{1, 2}$.

\begin{keydef}
An element $a_{ij}$ is said to have row index of $i$ and
column index of $j$.
\end{keydef}

As you probably have guessed, we assign special names to some
matrices depending on $m$ and $n$.

\begin{keydef}
\textbf{Square matrices} are matrices with equal number of rows
and columns. The elements with same row and column index are
called \textbf{diagonal element} of the matrix. The collection
of such elements is called \textbf{main diagonal}, or
\textbf{principal diagonal}.
\end{keydef}

What if either $m$ or $n$ is equal to $1$? Yes, we do have special
terms for them.

\begin{keydef}
\textbf{Row matrices} are matrices that is composed of only one
row. A \textbf{column matrix} is an array composed with one column.
\end{keydef}

If you have machine learning background, you likely have heard
of the term ``high dimensional arrays or vectors.'' Here, the
term ``dimension'' works the same way.

\begin{keydef}
The \textbf{dimension} of a row or column matrix is the number
of \textbf{components}, or the elements. A \textbf{$\mathbf{n}$-tuple}
is an $n$-dimensional row or column vector.
\end{keydef}

Let's take a look at an example from the previous matrix $B$
reproduced below.

\begin{minipage}{0.2\textwidth}
\[
B =
\begin{bNiceMatrix}[last-row,last-col]
\pi & \Vbrace[horizontal-label]{4}{4 \text{ rows}} \\
e \\
0 \\
1 \\
\Hbrace{1}{1 \text{ column}}
\end{bNiceMatrix}
\]
\end{minipage}
\hfill
\begin{minipage}{0.7\textwidth}
Here, matrix $B$ is a 4-dimensional matrix with four components.
This column matrix can also be called as a $4$-tuple.
\end{minipage}

We could also define matrices by its elements.

\begin{keydef}
A \textbf{zero matrix} is a matrix with all its elements as $0$.
\end{keydef}

With the fundamental terms in mind, let's discuss the properties
of matrices.

\section{Fundamental Properties}

First, when are two matrices ``equal'' to each other? Consider
the following definition for a formal explanation.

\begin{theobox}{Equality of Matrix}{Equality of Matrix}
Two matrices $A$ and $B$ with orders $m \times n$ and $p \times q$
respectively are said to be equal to each other if they retain
the same order and corresponding elements, i.e. they satisfy the
following properties.
\begin{enumerate}
\item $m = p$
\item $n = q$
\item
For elements $a_{ij}$ in $A$ and $b_{ij}$ in $B$, the following
equation satisfies.
\[
a_{ij} = b_{ij} \,
    \forall \, i \in [1, m] \text{ and } j \in [1, n]
\]
\end{enumerate}
\end{theobox}

Although this is very trivial and intuitive, it is always nice
to formally define concepts!

\begin{theobox}{Matrix Addition}{Matrix Addition}
For matrices $A$ and $B$ with same order, the element in the
$(i, j)$-entry of their sum $A + B$ is equivalent to the sum of
elements in the $(i, j)$-entry of $A$ and $B$.
\end{theobox}

Note that the orders of matrices being added must be identical.
We do not consider additions with matrices of different orders.
Below is an example of addition.
\[
\begin{bmatrix}
    1 & 3   & 5 \\
    2 & \pi & 4
\end{bmatrix}
+
\begin{bmatrix}
    -2 & 1   & 3 \\
    3  & \pi & -4
\end{bmatrix}
=
\begin{bmatrix}
    -1 & 4    & 8 \\
    5  & 2\pi & 0
\end{bmatrix}
\]
Similarly, we could see that subtractions are defined in similar
way. Subtracting the matrices above, we get the following results.
\[
\begin{bmatrix}
    1 & 3   & 5 \\
    2 & \pi & 4
\end{bmatrix}
-
\begin{bmatrix}
    -2 & 1   & 3 \\
    3  & \pi & -4
\end{bmatrix}
=
\begin{bmatrix}
    3  & 2 & 2 \\
    -1 & 0 & 8
\end{bmatrix}
\]
Now, as we look at the additions and subtractions of matrices,
you probably wondered about application of additive properties.
Intuitively, the fundamental properties of addition can be
applied with matrices addition.

\begin{theobox}{Commutative Property of Matrix Addition}
{Commutative Property of Matrix Addition}
For matrices $A$ and $B$ with order $m \times n$, $A + B = B + A$.
\end{theobox}
\begin{proof}
Let $A = [a_{ij}]_{m \times n}$ and $B = [b_{ij}]_{m \times n}$.
By definition, the sum $A + B = [a_{ij} + b_{ij}]$. Because
addition is commutative, the following equation holds.
\[
A + B
= [a_{ij} + b_{ij}]
= [b_{ij} + a_{ij}]
= B + A
\]
Thus, the commutative property of matrix addition holds.
\end{proof}

In addition to the commutative property, the associative property
is also applicable.

\begin{theobox}{Associative Property of Matrix Addition}
{Associative Property of Matrix Addition}
For matrices $A$, $B$, and $C$ with order $m \times n$,
$A + (B + C) = (A + B) + C$.
\end{theobox}
\begin{proof}
Let $A = [a_{ij}]_{m \times n}$, $B = [b_{ij}]_{m \times n}$,
and $C = [c_{ij}]_{m \times n}$. Consider the following equation.
\begin{align*}
    A + (B + C)
    &= [a_{ij}] + [b_{ij} + c_{ij}]
     = [a_{ij} + (b_{ij} + c_{ij})]
     = [a_{ij} + b_{ij} + c_{ij}] \\
    &= [(a_{ij} + b_{ij}) + c_{ij}]
     = [a_{ij} + b_{ij}] + [c_{ij}] \\
    &= (A + B) + C
\end{align*}
Thus, the associative property of addition applies to the matrix
addition.
\end{proof}

One obvious property in addition is that when you add $0$ to a number,
you get the number itself. The same applies in matrix addition.

\begin{theobox}{Additive Identity of Matrix Addition}
{Additive Identity of Matrix Addition}
The sum of matrices $A_{m \times n}$ and $0_{m \times n}$ yields
the matrix $A$, i.e. zero matrix is the identity elements of
matrix addition.
\end{theobox}
\begin{proof}
Let $A = [a_{ij}]_{m \times n}$. By definition, the sum $A + 0
= [a_{ij} + 0] = [a_{ij}] = A$. Therefore, zero matrix is the
identity element of matrix additions.
\end{proof}

Now what if we want to add same matrix multiple times? Fortunately,
we already resolved the issue. For instance, we write $x + x + x$
as $3x$. We can multiply matrices with scalars.

\begin{theobox}{Scalar Multiplication}{Scalar Multiplication}
For a complex number $\lambda$ and $A = [a_{ij}]$, $\lambda A
= [\lambda a_{ij}]$.
\end{theobox}

Let's take a look at an example.

\begin{probbox}{}{}
Compute $2A - 3B$ for $A$ and $B$ defined as the following.
\[
A =
\begin{bmatrix}
    2 & 3 & -1 \\
    1 & 7 & 3
\end{bmatrix},
\quad
B =
\begin{bmatrix}
    -2 & 1 & 0 \\
    5  & 3 & 8
\end{bmatrix}
\]
\end{probbox}
\begin{solution}
Before subtracting, we could perform scalar multiplication first.
\[
2A - 3B =
2\begin{bmatrix}
    2 & 3 & -1 \\
    1 & 7 & 3
\end{bmatrix}
-
3\begin{bmatrix}
    -2 & 1 & 0 \\
    5  & 3 & 8
\end{bmatrix}
=
\begin{bmatrix}
    4 & 6  & -2 \\
    2 & 14 & 6
\end{bmatrix}
-
\begin{bmatrix}
    -6 & 3 & 0 \\
    15 & 9 & 24
\end{bmatrix}
\]
From here we could subtract two matrices. Therefore, $2A - 3B$
can be represented as following.
\[
2A - 3B =
\begin{bmatrix}
    10  & 3  & -2 \\
    -13 & 5 & -18
\end{bmatrix}
\]
\end{solution}

With the scalar multiplication in mind, we can expand the application
of foundational concepts here with matrix.

\begin{theobox}{Additional Properties}{Additional Properties}
For $A = [a_{ij}]_{m \times n}$, $B = [b_{ij}]_{m \times n}$,
and $\lambda_1, \lambda_2 \in \mathbb{C}$, the following
properties hold.
\begin{enumerate}
\item $\lambda_1 (A \pm B) = \lambda_1 A \pm \lambda_1 B$
\item $(\lambda_1 \pm \lambda_2) A = \lambda_1 A \pm \lambda_2 A$
\item $\lambda_1 (\lambda_2 A) = (\lambda_1 \lambda_2) A$
\end{enumerate}
\end{theobox}

Let's prove them! Below is the proof for the first property.

\begin{proof}
Consider the following equation.
\[
\lambda_1 (A \pm B)
= \lambda_1 ([a_{ij}] \pm [b_{ij}])
= \lambda_1 ([a_{ij} \pm b_{ij}])
\]
Note that individual elements are multiplied in scalar multiplication.
Moreover, because $\lambda_1 (a_{ij} \pm b_{ij}) = \lambda_1 a_{ij}
\pm \lambda_1 b_{ij}$, the following equation satisfy.
\[
\lambda_1 ([a_{ij} \pm b_{ij}])
= [\lambda_1 a_{ij} \pm \lambda_1 b_{ij}]
= [\lambda_1 a_{ij}] \pm [\lambda_1 b_{ij}]
= \lambda_1 A \pm \lambda_1 B
\]
Therefore, $\lambda_1 (A \pm B) = \lambda_1 A \pm \lambda_1 B$.
\end{proof}

Now let's prove the second property.

\begin{proof}
By definition of scalar multiplication, $(\lambda_1 \pm \lambda_2)
A$ can be represented as the following.
\[
(\lambda_1 \pm \lambda_2) A
= (\lambda_1 \pm \lambda_2) [a_{ij}]
= [(\lambda_1 \pm \lambda_2) a_{ij}]
\]
Using the distributive property, the following equation holds.
\[
[(\lambda_1 \pm \lambda_2) a_{ij}]
= [\lambda_1 a_{ij} \pm \lambda_2 a_{ij}]
= [\lambda_1 a_{ij}] \pm [\lambda_2 a_{ij}]
= \lambda_1 A \pm \lambda_2 A
\]
Thus, the equation $(\lambda_1 \pm \lambda_2) A =
\lambda_1 A \pm \lambda_2 A$ holds.
\end{proof}

Finally, here is the prove to the third property.

\begin{proof}
Consider the following equation.
\[
\lambda_1 (\lambda_2 [a_{ij}])
= \lambda_1 [\lambda_2 a_{ij}]
= [\lambda_1 \lambda_2 a_{ij}]
= [(\lambda_1 \lambda_2) a_{ij}]
= (\lambda_1 \lambda_2) [a_{ij}]
= (\lambda_1 \lambda_2) A
\]
Therefore, the property holds.
\end{proof}

We have discussed fundamental properties so far. In the next
section of the note, let's discuss what matrix multiplication
with intuitions and computations.







\end{document}


